\subsection{Проектирование баз данных}
\label{sec:design:database}

Проектирование базы данных выполнено в два этапа: концептуальное и логическое. На концептуальном этапе построена \textit{ER}-модель в нотации Питера Чена. Модель описывает предметную область в терминах сущностей, их атрибутов и связей и не зависит от конкретной СУБД. На логическом этапе модель переведена в набор реляционных таблиц с первичными и внешними ключами [4]. 
Физическая реализация схемы в \textit{PostgreSQL} рассмотрена в подразделе~\ref{sec:development:database}.

Информационное пространство системы делится на три группы данных. Первая описывает пользователей и то, что они отслеживают. Вторая накапливается при работе системы: объявления с площадки, найденные совпадения и журнал отправленных уведомлений. К третьей относится справочник характеристик категорий площадки, который система собирает сама и по которому строит меню фильтров. На концептуальном уровне все три группы входят в одну модель, поскольку связаны между собой через отслеживаемый товар.

Центральная сущность модели~-- «Пользователь». С ней связана сущность «Отслеживаемый товар», которая хранит критерии поиска. Сущность «Объявление» содержит сведения, полученные с площадки, и существует независимо от пользователей: одно объявление может подойти под товары нескольких людей. Связь товара и объявления вынесена в отдельную сущность «Совпадение», поскольку у самой связи есть собственные атрибуты: тип события, показатель выгоды, состояние доставки. Факт отправки сообщения фиксирует сущность «Уведомление». Справочную часть составляют сущности «Атрибут фильтра» и «Значение атрибута». \textit{ER}-диаграмма приведена на рисунке~\ref{fig:design:er}. На концептуальном уровне сущности, связи и атрибуты названы по-русски и от имён таблиц базы данных не зависят; соответствие имён приведено ниже, при описании логической схемы.

\begin{figure}[H]
	\centering
	\includegraphics[width=\linewidth]{kufar/er_chen}
	\caption{\textit{ER}-диаграмма базы данных в нотации Питера Чена}
	\label{fig:design:er}
\end{figure}

Сущность «Пользователь»
\begin{itemize}
	\item идентификатор~-- первичный ключ, внутренний номер пользователя в системе;
	\item идентификатор \textit{Telegram}~-- номер учётной записи, по которому бот узнаёт пользователя; должен быть уникальным;
	\item имя учётной записи~-- публичное имя в \textit{Telegram} (может отсутствовать);
	\item имя из профиля~-- имя, которое бот использует в обращениях;
	\item язык интерфейса~-- язык меню и уведомлений (русский или английский);
	\item признак активности~-- снимается, если пользователь заблокировал бота, чтобы система не пыталась отправлять ему сообщения;
	\item признак уведомлений~-- включена ли отправка уведомлений в \textit{Telegram};
	\item электронная почта~-- адрес для копий уведомлений и отчётов (может отсутствовать);
	\item признак подтверждения почты~-- подтверждён ли адрес кодом; без подтверждения письма на адрес не отправляются;
	\item признак почтовых копий~-- отправлять ли копии уведомлений на подтверждённый адрес;
	\item хеш кода подтверждения~-- \textit{HMAC}-свёртка шестизначного кода; сам код в базе данных не хранится;
	\item срок действия кода~-- момент, после которого код подтверждения становится недействительным;
	\item время отправки кода~-- момент последней отправки; ограничивает частоту повторных запросов;
	\item число попыток ввода кода~-- после пятой неудачной попытки код нужно запросить заново;
	\item дата регистрации~-- момент первого обращения пользователя к боту;
	\item дата изменения~-- момент последнего изменения записи.
\end{itemize}

Сущность «Отслеживаемый товар»
\begin{itemize}
	\item идентификатор~-- первичный ключ, уникальный номер товара;
	\item владелец~-- ссылка на пользователя, который добавил товар;
	\item название~-- название товара, по которому ведётся текстовый отбор; в категориях с фильтрами площадки может быть пустым;
	\item категория~-- одна из категорий фиксированного перечня (велосипеды, ноутбуки, аренда квартир, электроника, инструменты и другие);
	\item пороговая цена~-- цена в белорусских рублях, ниже которой объявление считается выгодным;
	\item слова поиска~-- список слов, каждое из которых должно встретиться в заголовке или описании объявления;
	\item слова-исключения~-- список слов, при наличии которых объявление отклоняется;
	\item фильтры площадки~-- набор кодов характеристик и выбранных значений (бренд, размер рамы, число комнат, диапазон площади), которые передаются площадке в поисковом запросе;
	\item признак активности~-- снимается, когда пользователь приостанавливает отслеживание;
	\item дата добавления~-- момент добавления товара; объявления, опубликованные позже этого момента, считаются новыми.
\end{itemize}

Сущность «Объявление»
\begin{itemize}
	\item идентификатор~-- первичный ключ, внутренний номер объявления;
	\item площадка~-- источник объявления; в текущей версии системы это \textit{Kufar.by};
	\item внешний идентификатор~-- номер объявления на площадке; вместе с площадкой однозначно определяет объявление;
	\item заголовок~-- заголовок объявления, заданный продавцом;
	\item описание~-- текст объявления, а при его отсутствии перечень характеристик товара;
	\item цена~-- цена в белорусских рублях (может отсутствовать, если цена договорная);
	\item валюта~-- валюта, в которой указана цена;
	\item ссылка~-- адрес страницы объявления на площадке;
	\item изображения~-- список адресов фотографий товара;
	\item местоположение~-- регион и город или район;
	\item сведения о продавце~-- имя и номер учётной записи продавца на площадке;
	\item состояние~-- состояние товара, указанное продавцом (новое, бывшее в употреблении);
	\item исходный ответ площадки~-- полный ответ по объявлению в формате \textit{JSON};
	\item дата публикации~-- момент размещения объявления на площадке;
	\item первое появление~-- момент, когда система впервые получила объявление;
	\item последнее появление~-- момент последнего получения объявления; по нему удаляются устаревшие записи.
\end{itemize}

Сущность «Совпадение»
\begin{itemize}
	\item идентификатор~-- первичный ключ, уникальный номер совпадения;
	\item отслеживаемый товар~-- ссылка на товар, критериям которого удовлетворило объявление;
	\item объявление~-- ссылка на подошедшее объявление;
	\item показатель выгоды~-- доля, на которую цена объявления ниже пороговой (от нуля до единицы);
	\item подробности отбора~-- тип события (новое объявление или снижение цены), прежняя цена и порог на момент отбора;
	\item время доставки~-- момент успешной отправки уведомления в \textit{Telegram}; пустое значение означает, что совпадение ждёт в очереди;
	\item число попыток~-- количество неудачных отправок; после пятой попытки совпадение исключается из очереди;
	\item текст ошибки~-- причина последней неудачной отправки;
	\item время почтовой копии~-- момент отправки копии уведомления на электронную почту;
	\item число почтовых попыток~-- счётчик неудачных попыток для отдельной почтовой очереди;
	\item оценка пользователя~-- оценка предложения кнопкой под уведомлением («Интересно», «Не то», «Куплю»);
	\item дата создания~-- момент, когда объявление прошло отбор.
\end{itemize}

Сущность «Уведомление»
\begin{itemize}
	\item идентификатор~-- первичный ключ, уникальный номер записи журнала;
	\item получатель~-- ссылка на пользователя, которому отправлено сообщение;
	\item отслеживаемый товар~-- ссылка на товар, по которому пришло уведомление;
	\item объявление~-- ссылка на объявление из уведомления;
	\item совпадение~-- ссылка на совпадение, по которому отправлено сообщение (может отсутствовать);
	\item площадка и внешний идентификатор~-- копия идентификации объявления, по которой проверяется, получал ли пользователь это объявление раньше;
	\item цена при отправке~-- цена объявления на момент отправки; повторно объявление отправляется только по более низкой цене;
	\item показатель выгоды~-- показатель выгоды совпадения на момент отправки;
	\item дата отправки~-- используется для часового и суточного ограничений и для удаления записей старше 90 дней.
\end{itemize}

Сущность «Атрибут фильтра»
\begin{itemize}
	\item идентификатор~-- первичный ключ, уникальный номер характеристики;
	\item категория~-- код категории \textit{Kufar.by}, к которой относится характеристика;
	\item внутреннее имя~-- имя характеристики на площадке; вместе с категорией должно быть уникальным;
	\item код запроса~-- короткий код, под которым площадка принимает характеристику в поисковом запросе;
	\item название~-- название характеристики для меню бота (например, «Производитель»);
	\item частота~-- число объявлений выборки, в которых встретилась характеристика;
	\item дата обновления~-- момент последнего сбора справочника.
\end{itemize}

Сущность «Значение атрибута»
\begin{itemize}
	\item идентификатор~-- первичный ключ, уникальный номер значения;
	\item атрибут~-- ссылка на характеристику, которой принадлежит значение;
	\item код значения~-- код, который передаётся площадке в запросе; вместе с атрибутом должен быть уникальным;
	\item название значения~-- название для кнопки меню (например, «\textit{Trek}»);
	\item родительский регион~-- код региона для значений «город» и «район» (может отсутствовать);
	\item частота~-- число объявлений выборки с этим значением; определяет порядок значений в меню;
	\item дата обновления~-- момент последнего сбора справочника.
\end{itemize}

Связи между сущностями имеют разную мощность. Сущности «Пользователь» и «Отслеживаемый товар» связаны отношением «один ко многим»: один пользователь отслеживает несколько товаров, а каждый товар принадлежит одному пользователю. Между сущностями «Отслеживаемый товар» и «Объявление» существует отношение «многие ко многим»: одно объявление может подойти под товары разных пользователей, а по одному товару находится много объявлений. Это отношение реализовано через сущность «Совпадение», которая связана с товаром и с объявлением отношениями «один ко многим». Для каждой пары товара и объявления хранится одно совпадение: при снижении цены система обновляет его, а не создаёт новое.

Сущности «Совпадение» и «Уведомление» связаны отношением «один ко многим». Если объявление подешевело после первой отправки, пользователь должен получить по тому же совпадению ещё одно уведомление с новой ценой, и в журнале появится вторая запись. В журнале уведомлений получатель хранится напрямую, чтобы часовое и суточное ограничения подсчитывались без обращения к таблицам товаров и совпадений. На концептуальной схеме эта связь не показана: пользователь однозначно определяется через совпадение и отслеживаемый товар.

В справочной части сущности «Атрибут фильтра» и «Значение атрибута» связаны отношением «один ко многим». Сущности «Значение атрибута» и «Отслеживаемый товар» образуют отношение «многие ко многим»: товар уточняется несколькими значениями, а одно значение (например, бренд) выбрано в товарах разных пользователей. В отдельную таблицу эта связь не вынесена. Выбранные коды значений хранятся в самом товаре как набор в формате \textit{JSON}, поскольку в таком же виде они подставляются в поисковый запрос к площадке, а справочник пересобирается каждую неделю и не должен каскадно менять товары пользователей.

Логическая схема базы данных приведена на рисунке~\ref{fig:design:logical}. Каждая сущность стала таблицей с суррогатным первичным ключом \textit{id}, а связи «один ко многим» выражены внешними ключами. Сущности «Пользователь», «Отслеживаемый товар», «Объявление», «Совпадение» и «Уведомление» соответствуют таблицам \textit{users}, \textit{tracked\_products}, \textit{listings}, \textit{matches} и \textit{sent\_notifications}, справочная часть~-- таблицам \textit{kufar\_attributes} и \textit{kufar\_attribute\_values}. Кроме таблиц, соответствующих сущностям модели, в схему входит резервная таблица пользовательских атрибутов товара \textit{custom\_attributes}. Она создана для будущего расширения системы, в текущей версии не заполняется и в концептуальную модель не включена.

\begin{figure}[H]
	\centering
	\includegraphics[width=\linewidth]{kufar/physical}
	\caption{Логическая схема базы данных}
	\label{fig:design:logical}
\end{figure}

Разделение данных на пользовательские, рабочие и справочные определяет дальнейшую работу со схемой. Пользовательские таблицы меняются только по командам пользователя и невелики. Рабочие таблицы растут с каждым циклом опроса, поэтому для них предусмотрено удаление устаревших записей и индексы под запросы цикла. Справочные таблицы перезаписываются по расписанию и на пользовательские данные не ссылаются.
